\section{Тестирование и результаты}
\label{sec:testing}

Тесты решают две задачи: подтвердить корректность и показать работающий результат под лимитом памяти.
Стек --- Gradle, JUnit~5, AssertJ. Всего $106$ тестов, все зелёные.

\subsection{Синтетические тесты (кратко)}
\begin{itemize}[itemsep=3pt]
  \item \textbf{Эталон.} Главный инструмент проверки --- небольшая заведомо верная плотная реализация
        LeaderRank: строится матрица перехода $(n{+}1)\times(n{+}1)$ и считается степенной метод.
        Сложность $\bigO(n^2)$ ограничивает её малыми графами, но на них она даёт <<истину>>, с
        которой сверяется основной движок. На примере \texttt{small.csv} расхождение
        $3{,}3\cdot10^{-15}$, то есть на уровне ошибки округления.
  \item \textbf{Графы вручную.} Треугольник, звезда (это и есть гипер-узел, центр обязан получить
        наибольший ранг), цепочка, висячая вершина, петля, две несвязные компоненты. Каждый случай
        фиксирует своё свойство.
  \item \textbf{Инварианты.} Сохранение массы ($\sum s\approx n$), положительность рангов, сходимость
        за ограниченное число итераций.
  \item \textbf{Главный тест --- инвариантность к числу потоков.} Запуск с $P=1$ и $P=8$ даёт байт в
        байт одинаковый результат. Один этот тест сразу доказывает и корректность параллелизма (нет
        гонок), и детерминизм (нет зависимости от порядка).
  \item \textbf{Генератор данных.} Свой детерминированный генератор R-MAT со степенным распределением
        --- чтобы в графе реально были гипер-узлы. Зерно фиксировано, граф повторяется.
  \item \textbf{Тест под лимитом памяти.} Полный конвейер запускается как отдельный процесс под
        лимитом cgroup. Тест проверяет завершение и пиковый RSS под лимитом. Если на машине нет cgroup
        или \texttt{systemd-run} не может выставить лимит, тест автоматически пропускается, а не
        падает.
\end{itemize}

\subsection{Реальные графы}
Замеры сделаны скриптом \texttt{scripts/benchmark.sh}. Стенд: $24$ ядра, $\approx 29$~ГБ RAM, SSD,
OpenJDK~21, Serial GC. Движок параллельный out-of-core, бюджет памяти берётся из \texttt{-Xmx}. Время
\texttt{compute} --- <<горячее>> (последний прогон при \texttt{-{}-repeat=5}). Время
\texttt{preprocess} однократное, включает разбор CSV и внешнюю сортировку.

\paragraph{Масштабирование по потокам.} Граф R-MAT (scale=20), $16$ млн рёбер, $n = 640\,324$. Число
итераций --- $38$ при любом $P$, результат бит в бит не зависит от $P$.

\begin{center}
\small
\begin{tabular}{rrrrrr}
\toprule
$P$ & preprocess, мс & compute, мс & ускорение & эффективность & пик RSS, МиБ\\
\midrule
1  & 3450{,}9 & 1709{,}7 & 1{,}00 & 100\% & 260{,}6\\
2  & 1911{,}8 & 915{,}5  & 1{,}87 & 93\%  & 220{,}0\\
4  & 1794{,}1 & 543{,}7  & 3{,}14 & 79\%  & 218{,}6\\
8  & 1754{,}6 & 371{,}0  & 4{,}61 & 58\%  & 219{,}7\\
16 & 1758{,}8 & 269{,}2  & 6{,}35 & 40\%  & 222{,}0\\
24 & 1772{,}3 & 231{,}4  & \textbf{7{,}39} & 31\% & 245{,}6\\
\bottomrule
\end{tabular}
\end{center}

Счётное ядро --- это потоковый SpMV, упирающийся в пропускную способность памяти. Поэтому
эффективность падает с ростом $P$: полоса памяти насыщается около $P=8$--$16$. Это честное свойство
задачи, а не баг. Препроцессинг ускоряется примерно вдвое и выходит на плато: остаток --- это
последовательная сборка id-карты и сортировка. Пик RSS почти не зависит от $P$.

\paragraph{Реальные графы SNAP.} $P=24$, \texttt{-Xmx2g}, лимит итераций $100$.

\begin{center}
\small
\begin{tabular}{lrrrrr}
\toprule
граф & $n$ & $m$ & preprocess, мс & compute, мс & пик RSS, МиБ\\
\midrule
web-Google        & $875\,713$   & $5\,105\,039$   & 703{,}9  & 405{,}3  & 237{,}9\\
soc-LiveJournal1  & $4\,847\,571$ & $68\,993\,773$ & 8174{,}7 & 3327{,}9 & 1235{,}1\\
\bottomrule
\end{tabular}
\end{center}

Оба графа считаются out-of-core: рёбра не держатся в памяти. soc-LiveJournal1 с $69$ млн рёбер
обрабатывается под кучей $2$~ГБ.

\paragraph{Сходимость на реальных графах.} Критерий остановки масштабируется числом вершин:
$\sum_i|s_{\mathrm{new}}[i] - s_{\mathrm{old}}[i]| < \mathit{tol}\cdot n$, $\mathit{tol}=10^{-8}$. Это
важно: масса сохраняется и равна $n$, поэтому ненормированная разность растёт с $n$. Фиксированный
абсолютный порог на большом графе требовал бы изменения меньше уровня шума округления и никогда бы не
срабатывал. Деление на $n$ делает <<сошлось>> одинаковым на любом масштабе. Истинное число итераций до
сходимости (замерено поднятием лимита):

\begin{center}
\begin{tabular}{lr}
\toprule
граф & итераций до сходимости\\
\midrule
R-MAT (scale=20)  & 38\\
web-Google        & 191\\
soc-LiveJournal1  & 1909\\
\bottomrule
\end{tabular}
\end{center}

Реальные графы, особенно социальная сеть, имеют меньший спектральный зазор, поэтому степенной метод
сходится медленнее. Оценка <<десятки итераций>> верна для графов вроде R-MAT, но не для соцсети. При
этом ранжирование лидеров стабилизируется намного раньше численной сходимости: top-5 лидеров
совпадает (те же вершины и тот же порядок) при лимите $100$ и при полной сходимости для обоих графов.
Медленный <<хвост>> лишь уточняет младшие разряды очков, не меняя порядок лидеров. Значит лимит $100$
даёт правильных лидеров даже на LiveJournal, а флаг \texttt{-{}-max-iterations} позволяет при желании
довести счёт до полной сходимости.

\paragraph{RSS в зависимости от \texttt{-Xmx}.} На том же R-MAT при сжатии кучи с $2$~ГБ до
$128$~МБ время \texttt{compute} остаётся постоянным ($\approx 215$--$226$~мс), а RSS держится в районе
$190$--$305$~МиБ. Движок не деградирует под давлением памяти, а RSS не следит за размером графа (рёбра
на диске). Главный рычаг RSS --- это \texttt{-Xmx}. Нижняя граница памяти: web-Google считается до
\texttt{-Xmx48m}, soc-LiveJournal1 до \texttt{-Xmx224m}. На границе всё упирается уже в честный
$\bigO(n)$ (векторы значений фазы счёта), а не в предобработку --- ровно тот предел задачи,
обозначенный в разделе~\ref{sec:limits}.
